\documentclass[../../main.tex]{subfiles}
\graphicspath{{../../image/}{../../figure/}}
\begin{document}

\section{注记}\label{section:ch17-notes}

补偿紧致性方法由 Murat~[1] 与 Tartar~[1,2] 引入。将该方法用于构造双曲守恒律解的纲领由 Tartar~[2,3] 设计，他奠定了基本条件\eqref{align:ch17-17310}，并在标量情形的框架下展示了其使用。首次对方程组的应用归功于 DiPerna~[6]，它为这些思想的深入发展提供了动力，并已产生大量研究文献。这里的叙述只是浅尝辄止。清晰的导引还可见于 Evans 的讲义~[1]、Hörmander 的教材~[2]、Malek、Nečas 与 Růžička 的专著~[1]，以及 M.E.~Taylor 的论著~[2]。欲了解该主题更细致、更深入的发展，读者可参阅 Serre 的书~[11] 与 Yun-guang~Lu 的专著~[1]。Gui-Qiang~Chen~[8] 的综述文章对该领域研究趋势作了富有信息的介绍。

Young 测度由 L.C.~Young~[1] 引入。此处的叙述沿袭 Ball~[3]，读者可在其中找到超越 $L^\infty$ 框架的推广，以及对其他构造方式的评注与文献。

欲了解补偿紧致性理论的导引，见 Tartar~[1,2,3] 的讲义。div-curl 引理归功于 Murat 与 Tartar。此处给出的证明取自 Evans~[1]。引理17.2.2通常被称为 Murat 引理（Murat~[2]）。

测度值解的概念由 DiPerna~[11] 引入。关于该理论的进一步发展及其在守恒律方程组（包括对相变进行混合类型建模的方程组）解构造中的应用，见 Chen 与 Frid~[2]、Coquel 与 LeFloch~[1]、Demengel 与 Serre~[1]、Frid~[3]、Poupaud 与 Rascle~[1]、Royerburg 与 Slemrod~[1]、Schochet~[2] 与 Szepessy~[1]。

标量守恒律最早由 Tartar~[2] 用补偿紧致性方法处理。\cref{theorem:ch17-标量补偿紧致}证明中所使用的巧妙论证，由 Tartar（1986年5月致作者的私人通信）与 Chen 与 Lu~[1] 独立发现。另见 Vecchi~[1]。该方法已由 Tadmor、Rascle 与 Bagnerini~[1] 推广到二维空间标量守恒律，为此需要使用两个熵--熵流对。Yang、Zhu 与 Zhao~[2] 在 $L^p$ 框架下处理标量守恒律。通量仅连续的情形见 Panov~[7]。

带奇异源项的标量平衡律见 Schonbek~[2]、Amadori 与 Coclite~[1] 与 Kwon~[2]。

多个空间维数中标量守恒律的 Cauchy 问题也可经由测度值解在 $L^\infty$ 中求解（DiPerna~[11]、Szepessy~[2]、Panov~[2]）。Hwang 与 Tzavaras~[1] 采用了一种把动理公式与补偿紧致性理论思想相结合的替代途径。初边值问题由 Szepessy~[1] 以及 Ben Moussa 与 Szepessy~[1] 考虑。

标量守恒律中黏性与色散之间的竞争，在一维情形由 Schonbek~[1] 研究，在多个空间维数情形由 Kondo 与 LeFloch~[2]、LeFloch 与 Natalini~[1] 以及 Hwang 与 Tzavaras~[1] 研究。Camassa--Holm 浅水波方程解奇异极限的相应问题由 Coclite 与 Karlsen~[1] 讨论。

标量守恒律被四阶微分算子扰动的其他奇异扰动情形，已由补偿紧致性方法处理，包括 Tadmor~[3] 关于消失超黏性的文章，以及 Otto 与 Westdickenberg~[1] 关于薄膜逼近的论文。

近年来对双曲守恒律松弛的活跃研究已产生卷帙浩繁的文献，因此在此不可能穷尽列举。Natalini~[3] 的综述论文包含详尽的文献目录。若干相关文献已录于\cref{section:ch05-初边值问题}与\cref{section:ch06-注记}。
该理论发展中的开创性作用由 Tai-Ping~Liu~[21] 的工作所扮演，其动机来自 Whitham~[2]。补偿紧致性方法在此语境下首先由 Gui-Qiang~Chen 与 Tai-Ping~Liu~[1] 以及 Chen、Levermore 与 Liu~[1] 用于松弛形式为标量守恒律的两个守恒律方程组。半线性方程组\eqref{align:ch17-1751}（对理论与计算目的而言）的特殊效力首先由 Jin 与 Xin~[1] 认识到。17.5节对该方程组的处理改编自 Chen、Levermore 与 Liu~[1]、Lattanzio 与 Marcati~[1,2] 以及 Coquel 与 Perthame~[1] 的分析。另见 Yun-guang~Lu~[4]。关于燃烧模型系统\eqref{align:ch07-燃烧模型}的应用，
见 Hanouzet、Natalini 与 Tesei~[1]；关于色谱方程的应用，见 Collet 与 Rascle~[1] 以及 Klingenberg 与 Lu~[1]。此外，Lu 与 Klingenberg~[1]、Tzavaras~[4,5]、Gosse 与 Tzavaras~[1]、Serre~[15] 以及 Lattanzio 与 Serre~[2] 把补偿紧致性方法应用于松弛形式为两个守恒律方程组的三个或四个守恒律方程组。

对松弛理论的有趣贡献还包括 Coquel 与 Perthame~[1]、Marcati 与 Natalini~[2]、Marcati 与 Rubino~[3]、Rubino~[3]、Luo 与 Xin~[1] 以及 Luo 与 Yang~[3]。

半线性方程组\eqref{align:ch17-1751}的 $L^1$-Lipschitz 估计\eqref{align:ch17-17515}（它导致在 $BV$ 空间框架下处理松弛问题）归功于 Natalini~[1]。上半平面上非线性方程组\eqref{align:ch05-relaxed-system}的 $BV$ 解存在性由 Dafermos~[25] 确立，
但允许过渡到 $\mu\downarrow0$ 松弛极限、且与 $\mu$ 无关的 $BV$ 估计，目前仅对特殊情形 $p(u)=-u^{-1}$（Luo、Natalini 与 Yang~[1]、Amadori 与 Guerra~[2]）已知。其他已在 $BV$ 中处理的特殊方程组见 Tveito 与 Winther~[2]、Gosse~[1] 以及 Luo 与 Natalini~[1]。

17.6节对两个守恒律真非线性方程组的处理，以及17.7节对带单个拐点的等熵弹性方程组的处理，沿袭 DiPerna~[8] 的开创性论文。另见 Gripenberg~[1] 以及 Chen、Li 与 Li~[1]。当 $\sigma''(u)$ 在 $u=0$ 处不连续时\cref{theorem:ch17-弹性补偿紧致}的反例由 Greenberg~[3] 以及 Greenberg 与 Rascle~[1] 给出。

等熵弹性方程组在 $L^p$ 框架下由 J.W.~Shearer~[1]、Peixiong~Lin~[1] 以及 Serre 与 Shearer~[1] 处理。$L^\infty$ 中解的一种替代、原创构造（Demoulini、Stuart 与 Tzavaras~[1]、Miroshnikov 与 Tzavaras~[1,2]）基于这样的观察：离散化时间变量所得的方程组可通过变分原理求解。$L^\infty$ 中的初边值问题由 Heidrich~[1] 解决。

经由最大值原理的不变区域理论归功于 Chueh、Conley 与 Smoller~[1]（另见 Hoff~[2]）。系统讨论及若干例子见 Serre~[11]。松弛格式的稳定性与不变区域存在性之间的联系在 Serre~[15] 中讨论。\cref{theorem:ch17-弹性不变区域}的证明取自 Dafermos~[16]。另见 Serre~[3] 与 Venttsel'~[1]。

等熵气体动力学方程组最早由 DiPerna~[9] 用补偿紧致性方法处理，针对的是绝热指数 $\gamma=1+\frac{2}{n}$、$n=2\ell+1$ 的特殊值。随后 Gui-Qiang~Chen~[1] 以及 Ding、Chen 与 Luo~[1] 把分析推广到区间 $(1,\frac{5}{3}]$ 内任意 $\gamma$。综述见 Gui-Qiang~Chen~[2]。$\gamma\geqslant3$ 的情形由 Lions、Perthame 与 Tadmor~[1] 解决，而如\cref{theorem:ch17-气体存在性}所述的全范围 $1<\gamma<\infty$ 由 Lions、Perthame 与 Souganidis~[1] 涵盖。等温情形 $\gamma=1$ 是奇异的，由 Huang 与 Wang~[2] 以及 LeFloch 与 Shelukhin~[1] 处理。这里给出的特殊情形 $\gamma=3$ 的论证由 Gui-Qiang~Chen 告知作者。Vasseur~[1] 展示了对该特殊 $\gamma$ 值的额外正则性。更一般的、对非多方气体的真非线性方程组\eqref{align:ch07-等熵Euler}，
在压力函数 $p(\rho)$ 与函数 $\kappa\rho^\gamma$（连同其至多到四阶的导数）当 $\rho\to0$ 时渐近重合的假设下，由 Chen 与 LeFloch~[2,3] 处理。

物理黏性（相对于人工黏性）的使用，其顶点即\cref{theorem:ch17-气体物理黏性}，较晚出并归功于 Chen 与 Perepelitsa~[1]。另见 Chen 与 Perepelitsa~[2] 关于浅水波方程。这些工作本质地使用了 LeFloch 与 Westdickenberg~[1] 为处理径向对称情形所引入的技术，从而导出\cref{theorem:ch17-径向气体}。同一问题上，Chen 与 Perepelitsa~[3] 处理波的聚焦问题，并表明尽管方程组在 $r=0$ 处奇异，原点处仍不形成质量集中。（与\cref{section:ch16-注记}比较。）

Serre~[2,11] 的方法已使补偿紧致性方法足够灵活，甚至可处理特征族线性退化、严格双曲性失效等情形下的两个守恒律方程组。许多有趣方程组的解由 Chen 与 Wang~[1,2]、Dehua~Wang~[1,2]、Chen 与 Kan~[1]、Pui-Tak~Kan~[1]、Kan、Santos 与 Xin~[1]、Heibig~[2]、Yun-guang~Lu~[1]、Marcati 与 Natalini~[2,3]、Marcati 与 Rubino~[1,2]、Rugbino~[1,2]、Huijiang~Zhao~[1]、Frid 与 Santos~[1,2] 以及 Yun-guang~Lu~[3] 构造。由于分析本质地依赖于丰富熵族的可获得性，该方法对多于两个守恒律方程组的应用，目前限于激波与稀疏波曲线（对至多两个特征族）重合的特殊方程组（Benzoni-Gavage 与 Serre~[1]），以及非常特殊状态方程的非等熵气体动力学方程组（Chen 与 Dafermos~[1]、Chen、Li 与 Li~[1]、Bereux、Bonnivet 与 LeFloch~[1] 以及 Frid、Holden 与 Karlsen~[1]）。

补偿紧致性方法还成功应用于定常无旋超声速-亚声速或跨声速气体流的 Euler 方程组。该情形下方程组退化或变型；见 Morawetz~[2,4]、Chen、Dafermos、Slemrod 与 Wang~[1]、Chen、Slemrod 与 Wang~[1]、Huang 与 Wang~[1] 以及 Chen、Huang 与 Wang~[1]。如我们将在\cref{section:ch18-等距浸入}所见，补偿紧致性方法在微分几何中也找到了有趣应用。

对多种方程组，初值为周期的或 Riemann 数据的 $L^1$ 扰动的解的大时间行为，由 Chen 与 Frid~[1,3,4,6] 通过把标度不变性与紧致性相结合而确立。另见 Frid~[6] 以及 Frid 与 Rendón~[1]。补偿紧致性方法已被用于证明：带摩擦阻尼的 Euler 方程解的大时间行为由多孔介质方程支配；见 Serre 与 Xiao~[1]、Huang 与 Pan~[1,2]、Pan 与 Zhao~[1]、Luo 与 Yang~[1] 以及 Lattanzio 与 Rubino~[1]。关于带松弛方程组解的大时间行为，见 Serre~[19]。

动理公式已在\cref{chapter:ch06-标量守恒律L1理论}中有效应用于多个空间维数的标量守恒律，
并已成功推广到一维空间中的某些守恒律方程组，包括等熵气体流的 Euler 方程组（Berthelin 与 Bouchut~[1,2,3]）以及等熵弹动力学方程组（Perthame 与 Tzavaras~[1]、Tzavaras~[6]）。详细讨论与完备的文献目录见 Perthame~[2] 的专著。把动理公式与补偿紧致性理论的技术相结合，可导出解的精细性质。特别地，对两个守恒律的严格双曲方程组，Tzavaras~[6] 给出了两特征场之间振荡耦合的显式公式。

关于双曲守恒律中非线性效应的有价值洞见，来自对解算子如何与高度振荡初值相互作用的研究，譬如 $U_{0\varepsilon}(x)=V(x,x/\varepsilon)$，其中 $V(x,\cdot)$ 为周期、$\varepsilon$ 为小正参数。当方程组为线性时，快速振荡沿特征线输运，其幅值不衰减。相反极端下，当方程组严格双曲且真非线性时，17.4节与17.6节的分析表明：当 $\varepsilon\to0$ 时，所得解族 $U_\varepsilon(x,t)$ 包含一些子列，它们强收敛于以 $\{U_{0\varepsilon}\}$ 的弱极限为初值的解，即对 $t>0$，解算子淬灭初值的高频振荡。研究某些特征族可能线性退化、严格双曲性失效等的中间情形是有趣的。在众多特例研究（参见 Bonnefille~[1]、Gui-Qiang~Chen~[3,4,5]、E~[1]、Heibig~[1]、Rascle~[1]、Cheverry、Gues 与 Métivier~[1,2]、Bourdarias、Gisclon 与 Junca~[3] 以及 Serre~[5,8]）之后，一种连贯的振荡传播理论似乎正在形成（Serre~[11]）。

存在一种发展完善的振荡传播理论，它基于弱非线性几何光学方法，导出双曲方程组在高频、小振幅振荡初值下解的渐近展开。继 Landau~[1]、Lighthill~[1] 与 Whitham~[1] 的开创性工作之后，已涌现出大量纯形式、半严格或严格的文献，处理单相或可能共振的多相等。例如见 Choquet-Bruhat~[1]、Hunter 与 Keller~[1,2]、Majda 与 Rosales~[1]、Majda、Rosales 与 Schonbek~[1]、Pego~[4]、Hunter~[1]、Joly、Métivier 与 Rauch~[1,3] 以及 Cheverry~[1]。值得注意的是，即使解中发展出激波后，渐近展开仍然有效；见 DiPerna 与 Majda~[1]、Schochet~[5] 与 Cheverry~[2]。综述见 Majda~[5]，系统阐述见 Serre~[11]。多个空间维数标量守恒律非线性几何光学的有效性由 Chen、Junca 与 Rascle~[1] 确立。最近对基于\cref{chapter:ch14-前沿追踪与标准riemann半群}所讨论的标准 Riemann 半群的守恒律方程组几何光学的处理，见 Chen、Xiang 与 Zhang~[1]。

\end{document}
